\documentclass{aomart}

\usepackage[T1]{fontenc}
\usepackage[nopatch=footnote]{microtype}
\usepackage{amssymb,mathtools}
\usepackage{aliascnt}
\usepackage{enumitem}
\usepackage[nameinlink,noabbrev]{cleveref}

\allowdisplaybreaks
\hyphenpenalty=10000
\exhyphenpenalty=10000
\tolerance=2000
\emergencystretch=3em

\makeatletter
\def\@setsubjects{%
  \let\@subjects\@empty
  \ifx\@primarysubjects\@empty\else
    \edef\@subjects{Primary:~\@primarysubjects}%
  \fi
  \ifx\@secondarysubjects\@empty\else
    \ifx\@primarysubjects\@empty
      \edef\@subjects{Secondary:~\@secondarysubjects}%
    \else
      \edef\@subjects{\@subjects, Secondary:~\@secondarysubjects}%
    \fi
  \fi
  \ifpdf
    \hypersetup{pdfsubject=\@subjects}%
  \fi
  \if@aom@oldkeywords\else
    \@footnotetext{AMS Classification:~\@subjects.}%
  \fi
}
\makeatother

\setlist[enumerate]{topsep=4pt,itemsep=2pt}
\setlist[itemize]{topsep=4pt,itemsep=2pt}

\newtheorem{theorem}{Theorem}[section]
\newaliascnt{proposition}{theorem}
\newtheorem{proposition}[proposition]{Proposition}
\aliascntresetthe{proposition}
\newaliascnt{lemma}{theorem}
\newtheorem{lemma}[lemma]{Lemma}
\aliascntresetthe{lemma}
\newaliascnt{corollary}{theorem}
\newtheorem{corollary}[corollary]{Corollary}
\aliascntresetthe{corollary}
\theoremstyle{definition}
\newaliascnt{definition}{theorem}
\newtheorem{definition}[definition]{Definition}
\aliascntresetthe{definition}
\newaliascnt{example}{theorem}
\newtheorem{example}[example]{Example}
\aliascntresetthe{example}
\newaliascnt{standardfact}{theorem}
\newtheorem{standardfact}[standardfact]{Standard Fact}
\aliascntresetthe{standardfact}
\theoremstyle{remark}
\newaliascnt{remark}{theorem}
\newtheorem{remark}[remark]{Remark}
\aliascntresetthe{remark}

\crefname{theorem}{theorem}{theorems}
\crefname{proposition}{proposition}{propositions}
\crefname{lemma}{lemma}{lemmas}
\crefname{corollary}{corollary}{corollaries}
\crefname{definition}{definition}{definitions}
\crefname{example}{example}{examples}
\crefname{standardfact}{standard fact}{standard facts}
\crefname{remark}{remark}{remarks}

\newcommand{\CC}{\mathbb C}
\newcommand{\NN}{\mathbb N}
\newcommand{\ZZ}{\mathbb Z}
\newcommand{\Bcal}{\mathcal B}
\newcommand{\Ecal}{\mathcal E}
\newcommand{\Pcal}{\mathcal P}
\newcommand{\Rcal}{\mathcal R}
\newcommand{\Zcal}{\mathcal Z}
\newcommand{\Dfun}{\mathsf D}
\newcommand{\Efun}{\mathsf E}
\newcommand{\Sfun}{\mathsf S}
\newcommand{\Tfun}{\mathsf T}
\newcommand{\Cfun}{\mathsf C}
\newcommand{\Thetafun}{\mathsf\Theta}
\newcommand{\Spec}{\operatorname{Spec}}
\newcommand{\ord}{\operatorname{ord}}
\newcommand{\type}{\operatorname{type}}
\newcommand{\Log}{\operatorname{Log}}
\newcommand{\Res}{\operatorname{Res}}
\newcommand{\dist}{\operatorname{dist}}
\newcommand{\sech}{\operatorname{sech}}
\newcommand{\supp}{\operatorname{supp}}
\newcommand{\id}{\mathrm{id}}

\title[Cyclic analytic functions]{Cyclic Analytic Functions, Reciprocal Families, and Zero Geometry}
\author[Komorebi P.]{Komorebi P.}
\subject{primary}{msc2020}{30D20}
\subject{secondary}{msc2020}{33C20}
\subject{secondary}{msc2020}{30D15}
\subject{secondary}{msc2020}{11M35}
\subject{secondary}{msc2020}{05A40}
\keyword{cyclic projection}
\keyword{exponential polynomial}
\keyword{reciprocal polynomial family}
\keyword{Hadamard product}
\keyword{zero asymptotics}
\keyword{Appell sequence}
\copyrightnote{}
\doinumber{}

\pdfstringdefDisableCommands{\let\unskip\relax}
\hypersetup{
  pdftitle={Cyclic Analytic Functions, Reciprocal Families, and Zero Geometry},
  pdfauthor={Komorebi P.},
  pdfsubject={Cyclic projections, reciprocal families, entire functions, and zero geometry}
}

\begin{document}
\begin{abstract}
Fix a positive integer $N$.  Cyclic projection separates an analytic function into its coefficient classes modulo $N$.  Applied to the exponential, it gives a basis for the equation $y^{(N)}=\lambda^Ny$.  A signed version gives a basis for $y^{(N)}=-y$ and the entire denominator
\[
\Dfun_N(z)=\sum_{n\geq 0}(-1)^n\frac{z^{Nn}}{(Nn)!}.
\]
We study its reciprocal powers and show that their coefficients depend polynomially on the reciprocal order.  We also study the logarithmic derivative and its higher Riccati equation.  Normalized exponential sections give congruential reciprocal functions and Appell polynomial systems.  A scaled reciprocal construction gives a second polynomial family.  Writing $\Dfun_N(z)=\Delta_N(z^N)$ reduces the zero problem to an entire function of order $1/N$.  We prove a genus zero Hadamard product, locate every sufficiently large zero, show that all but finitely many zeros are simple, obtain the exact zero counting asymptotic, and derive meromorphic continuation of the symmetry grouped zero zeta.  The reciprocal coefficients are expressed through power sums of the reduced zeros.
\end{abstract}
\maketitle
\thispagestyle{plain}

\section{Introduction}\label{sec:introduction}

For a power series
\[
F(z)=\sum_{m\geq 0}a_mz^m
\]
and a positive integer $N$, the coefficients may be separated according to the residue of $m$ modulo $N$.  Root of unity averaging gives exact projection operators for this separation.  This elementary construction is usually called a multisection or a root filter.  It appears throughout the classical theory of generating functions and special functions.  Background on multisections, hypergeometric functions, and related coefficient systems may be found in \cite{Rainville1960,AndrewsAskeyRoy1999,Comtet1974}.

This article follows the consequences of applying the cyclic projections to the exponential function.  The projected exponentials form a finite differential system.  After one sign is inserted around the derivative cycle, they form a basis for
\[
y^{(N)}=-y.
\]
The section with initial value one is
\[
\Dfun_N(z)=\sum_{n\geq 0}(-1)^n\frac{z^{Nn}}{(Nn)!}.
\]
Its reciprocal powers produce polynomial coefficient families.  Its logarithmic derivative produces a higher Riccati equation.  The same reciprocal construction applied to every exponential section gives normalized congruential functions and ordinary Appell polynomial sequences.  Appell and binomial type sequences are discussed in \cite{Roman1984,Comtet1974}.

The second half of the paper concerns the zeros of $\Dfun_N$.  The finite exponential sum
\[
\Dfun_N(z)=\frac1N\sum_{\eta^N=-1}e^{\eta z}
\]
has a regular polygon of exponents.  Two adjacent exponentials dominate on each of $N$ rays.  Their cancellation gives one zero string on every ray.  We prove that these strings contain every sufficiently large zero and that every zero in them is simple.  This gives an exact linear zero counting law.

The rotational symmetry is more useful after the change of variable $w=z^N$.  There is an entire function
\[
\Delta_N(w)=\sum_{n\geq 0}(-1)^n\frac{w^n}{(Nn)!}
\]
with
\[
\Dfun_N(z)=\Delta_N(z^N).
\]
For $N\geq2$, the order of $\Delta_N$ is $1/N$.  Its Hadamard product therefore has genus zero.  The logarithm of this product identifies the transfer tangent coefficients with power sums of the reduced zeros.  The large zero formula also compares the reduced zero zeta with a Hurwitz zeta function.  Standard material on entire functions, canonical products, and exponential sums appears in \cite{Boas1954,Levin1980,Titchmarsh1939}.  We use singularity analysis in the form described in \cite{FlajoletSedgewick2009} and the classical continuation of the Hurwitz zeta function from \cite{Apostol1976}.

The article does not assign historical priority to every named coefficient family.  Several low order cases are classical, including the exponential, cosine, secant, tangent, Euler, Bernoulli type, and Appell cases.  The point is to place the full cyclic family in one analytic structure and to prove the relations among its differential equations, reciprocal powers, zeros, and coefficient systems.



\section{A map of the construction}\label{sec:map}

Two variables appear throughout the paper.  The variable $z$ carries the differential equation and the rotational symmetry.  The reduced variable
\[
w=z^N
\]
records one representative from each rotational orbit.  Confusing these variables hides the main simplification.  The denominator $\Dfun_N$ has order one as a function of $z$, while $\Delta_N$ has order $1/N$ as a function of $w$.

The analytic objects are introduced in the following order.  First,
\[
\Efun_{N,j}(z)
=
\sum_{n\geq0}\frac{z^{Nn+j}}{(Nn+j)!}
\]
is the residue $j$ part of the exponential.  Second,
\[
\Dfun_{N,j}(z)
=
\sum_{n\geq0}(-1)^n\frac{z^{Nn+j}}{(Nn+j)!}
\]
is the signed basis for $y^{(N)}=-y$.  The section $\Dfun_N=\Dfun_{N,0}$ is the denominator used in the reciprocal family
\[
\Sfun_N^{(r)}=\Dfun_N^{-r}.
\]
Its logarithmic derivative is $\Tfun_N$.

A different reciprocal family comes from normalizing each exponential section at the origin.  The function
\[
\Cfun_{N,j}(z)
=
\frac{z^j/j!}{\Efun_{N,j}(z)}
\]
is a unit series in $z^N$.  Its powers, after multiplication by $e^{xz}$, give Appell polynomials.  The theta family couples a scaled argument with a reciprocal order.

The zero theory begins only after these algebraic relations have been established.  The factorization
\[
\Dfun_N(z)=\Delta_N(z^N)
\]
turns every orbit of $N$ zeros in the $z$ plane into one zero in the $w$ plane.  The large zeros in the $z$ plane are located by two dominant exponentials.  The zeros in the $w$ plane then give a genus zero product.  The logarithm of that product returns to the coefficient theory through
\[
\text{zero power sums}
=
\text{logarithmic coefficients}.
\]
This equality is the bridge between the first and second halves of the paper.

All powers with a complex parameter are first defined as formal series at the origin.  An analytic interpretation requires a simply connected zero free region and one fixed logarithm.  Integer reciprocal powers are globally meromorphic and do not require a branch choice.

\section{Cyclic projection of analytic functions}\label{sec:cyclic_projection}

Fix $N\geq1$ and write
\[
\omega=e^{2\pi i/N}.
\]
Let $\Omega\subseteq\CC$ be invariant under multiplication by $\omega$.  For a holomorphic function $F$ on $\Omega$, define
\[
(\Bcal_{N,j}F)(z)
=
\frac1N\sum_{\ell=0}^{N-1}\omega^{-j\ell}F(\omega^\ell z),
\qquad 0\leq j<N.
\]
The same formula is valid for formal power series.

\begin{theorem}[Root filter]\label{thm:root_filter}
If
\[
F(z)=\sum_{m\geq0}a_mz^m,
\]
then
\[
\Bcal_{N,j}F(z)
=
\sum_{m\equiv j\bmod N}a_mz^m.
\]
The operators satisfy
\[
\Bcal_{N,i}\Bcal_{N,j}=\delta_{ij}\Bcal_{N,j}
\]
and
\[
\sum_{j=0}^{N-1}\Bcal_{N,j}=I.
\]
\end{theorem}

\begin{proof}
Substitution of the power series gives
\[
\Bcal_{N,j}F(z)
=
\sum_{m\geq0}a_mz^m
\left(
\frac1N\sum_{\ell=0}^{N-1}\omega^{(m-j)\ell}
\right).
\]
The inner average is one when $m\equiv j\bmod N$ and zero otherwise.  The operator identities follow from the disjoint residue classes.
\end{proof}

\begin{proposition}[Covariance characterization]\label{prop:covariance}
A holomorphic function $F$ belongs to the range of $\Bcal_{N,j}$ if and only if
\[
F(\omega z)=\omega^jF(z).
\]
\end{proposition}

\begin{proof}
Every monomial selected by $\Bcal_{N,j}$ has the stated covariance.  Conversely, write $F=\sum_k\Bcal_{N,k}F$.  Replacing $z$ by $\omega z$ gives
\[
\sum_k\omega^k\Bcal_{N,k}F
=
\omega^j\sum_k\Bcal_{N,k}F.
\]
Application of $\Bcal_{N,k}$ shows that $\Bcal_{N,k}F=0$ for $k\ne j$.
\end{proof}

Multiplication is not diagonal in these coordinates.  The packet labels add modulo $N$.

\begin{theorem}[Cyclic product law]\label{thm:cyclic_product}
If $F$ and $G$ are holomorphic on the same rotation invariant domain, then
\[
\Bcal_{N,k}(FG)
=
\sum_{i+j\equiv k\bmod N}
(\Bcal_{N,i}F)(\Bcal_{N,j}G).
\]
\end{theorem}

\begin{proof}
A product of a monomial of degree congruent to $i$ and a monomial of degree congruent to $j$ has degree congruent to $i+j$.  Summing the corresponding products gives the formula.
\end{proof}

\begin{remark}\label{rem:formal_analytic}
Every identity in this section holds formally.  In the holomorphic setting, the finite projector sums preserve local uniform convergence.  No choice of branch is involved.
\end{remark}

\section{The transfer exponential basis}\label{sec:transfer_exponential}

\begin{definition}\label{def:transfer_exponential}
For $\lambda\in\CC$ and $0\leq j<N$, define
\[
\Efun_{N,j}(\lambda z)
=
\Bcal_{N,j}(e^{\lambda z})
=
\sum_{n\geq0}\frac{(\lambda z)^{Nn+j}}{(Nn+j)!}.
\]
\end{definition}

\begin{theorem}[Differential cycle and initial data]\label{thm:exponential_cycle}
For $1\leq j<N$,
\[
\frac{d}{dz}\Efun_{N,j}(\lambda z)
=
\lambda\Efun_{N,j-1}(\lambda z),
\]
while
\[
\frac{d}{dz}\Efun_{N,0}(\lambda z)
=
\lambda\Efun_{N,N-1}(\lambda z).
\]
Thus
\[
\frac{d^N}{dz^N}\Efun_{N,j}(\lambda z)
=
\lambda^N\Efun_{N,j}(\lambda z).
\]
For $0\leq k<N$,
\[
\left.\frac{d^k}{dz^k}\Efun_{N,j}(\lambda z)\right|_{z=0}
=
\delta_{kj}\lambda^j.
\]
If $\lambda\ne0$, the $N$ functions $\Efun_{N,j}(\lambda z)$ form a basis of the solution space of $y^{(N)}=\lambda^Ny$.
\end{theorem}

\begin{proof}
Termwise differentiation gives the cyclic derivative formulas.  One full cycle contributes $\lambda^N$.  The first nonzero term of $\Efun_{N,j}(\lambda z)$ is $(\lambda z)^j/j!$, which gives the initial data.  The initial value vectors are linearly independent when $\lambda\ne0$.  The solution space of an order $N$ constant coefficient equation has dimension $N$, so these functions form a basis.
\end{proof}


\begin{corollary}[Wronskian of the transfer basis]\label{cor:wronskians}
For $\lambda\ne0$, the Wronskian of the ordered basis
\[
\Efun_{N,0}(\lambda z),\ldots,\Efun_{N,N-1}(\lambda z)
\]
is the constant
\[
\lambda^{N(N-1)/2}.
\]
\end{corollary}

\begin{proof}
At the origin, the derivative matrix is diagonal with entries
\[
1,\lambda,\ldots,\lambda^{N-1}.
\]
Abel's identity shows that the Wronskian is constant because the coefficient of $y^{(N-1)}$ in the differential equation is zero.
\end{proof}

\begin{corollary}[Addition law]\label{cor:exponential_addition}
For $0\leq k<N$,
\[
\Efun_{N,k}(\lambda(z+w))
=
\sum_{i+j\equiv k\bmod N}
\Efun_{N,i}(\lambda z)\Efun_{N,j}(\lambda w).
\]
\end{corollary}

\begin{proof}
Apply the diagonal cyclic projector to
\[
e^{\lambda(z+w)}=e^{\lambda z}e^{\lambda w}.
\]
The cyclic product law gives the right hand side.
\end{proof}

The basis also has a scalar hypergeometric form.

\begin{theorem}[Hypergeometric form]\label{thm:hypergeometric_form}
Let
\[
A_{N,j}
=
\left\{
\frac{j+m}{N}
\mathrel{:}
1\leq m\leq N
\text{ and }
j+m\ne N
\right\}.
\]
Then
\[
\Efun_{N,j}(\lambda z)
=
\frac{(\lambda z)^j}{j!}
{}_0F_{N-1}
\left(
\begin{matrix}-\\ A_{N,j}\end{matrix}
\mathrel{,}
\frac{(\lambda z)^N}{N^N}
\right).
\]
\end{theorem}

\begin{proof}
The multiplication formula for the gamma function gives
\[
(Nn+j)!
=
j!N^{Nn}n!
\prod_{a\in A_{N,j}}(a)_n.
\]
Insertion into the defining series gives the stated hypergeometric expansion.
\end{proof}

\begin{example}\label{ex:classical_exponential}
For $N=1$, the only section is $e^{\lambda z}$.  For $N=2$,
\[
\Efun_{2,0}(z)=\cosh z,
\qquad
\Efun_{2,1}(z)=\sinh z.
\]
These are classical instances of the same cyclic basis.
\end{example}

\section{The signed system and the universal denominator}\label{sec:signed_system}

\begin{definition}\label{def:signed_sections}
For $0\leq j<N$, define
\[
\Dfun_{N,j}(z)
=
\sum_{n\geq0}(-1)^n\frac{z^{Nn+j}}{(Nn+j)!}.
\]
The universal denominator is
\[
\Dfun_N(z)=\Dfun_{N,0}(z).
\]
\end{definition}

Fix any $\rho\in\CC$ with $\rho^N=-1$.

\begin{proposition}[Relation with the transfer exponential]\label{prop:signed_exponential_relation}
For every $j$,
\[
\Dfun_{N,j}(z)
=
\rho^{-j}\Efun_{N,j}(\rho z).
\]
\end{proposition}

\begin{proof}
The coefficient of $z^{Nn+j}$ on the right is
\[
\rho^{-j}\rho^{Nn+j}\frac1{(Nn+j)!}
=
(-1)^n\frac1{(Nn+j)!}.
\]
\end{proof}

\begin{theorem}[Signed derivative cycle]\label{thm:signed_cycle}
For $1\leq j<N$,
\[
\Dfun_{N,j}'=\Dfun_{N,j-1},
\]
and
\[
\Dfun_{N,0}'=-\Dfun_{N,N-1}.
\]
Every signed section satisfies
\[
\Dfun_{N,j}^{(N)}=-\Dfun_{N,j}.
\]
The $N$ signed sections form the basis of $y^{(N)}=-y$ determined by
\[
\Dfun_{N,j}^{(k)}(0)=\delta_{jk}
\]
for $0\leq j,k<N$.
\end{theorem}

\begin{proof}
The derivative formulas follow directly from the series.  The first nonzero term of $\Dfun_{N,j}$ is $z^j/j!$, so the initial value matrix is the identity.  The basis statement follows from the dimension of the solution space.
\end{proof}

\begin{corollary}[Wronskian of the signed basis]\label{cor:signed_wronskian}
The Wronskian of
\[
\Dfun_{N,0}(z),\ldots,\Dfun_{N,N-1}(z)
\]
is one.
\end{corollary}

\begin{proof}
The derivative matrix at the origin is the identity.  Abel's identity shows that the Wronskian is constant because the coefficient of $y^{(N-1)}$ in the differential equation is zero.
\end{proof}


\begin{proposition}[Twisted addition law]\label{prop:signed_addition}
For $0\leq k<N$,
\[
\Dfun_{N,k}(z+w)
=
\sum_{\substack{0\leq i,j<N\\ i+j\equiv k\bmod N}}
(-1)^{\lfloor(i+j)/N\rfloor}
\Dfun_{N,i}(z)\Dfun_{N,j}(w).
\]
\end{proposition}

\begin{proof}
Use \cref{prop:signed_exponential_relation,cor:exponential_addition}.  The factor introduced by a pair $(i,j)$ is
\[
\rho^{i+j-k}.
\]
Since $0\leq i+j\leq2N-2$ and $i+j\equiv k\bmod N$, this factor is one when $i+j=k$ and $-1$ when $i+j=k+N$.
\end{proof}

\begin{theorem}[Finite exponential representation]\label{thm:finite_exponential_representation}
Let
\[
\Rcal_N=\{\eta\in\CC\mathrel{:}\eta^N=-1\}.
\]
Then
\[
\Dfun_N(z)
=
\frac1N\sum_{\eta\in\Rcal_N}e^{\eta z}.
\]
More generally,
\[
\Dfun_{N,j}(z)
=
\frac{\rho^{-j}}N
\sum_{\ell=0}^{N-1}
\omega^{-j\ell}e^{\rho\omega^\ell z}.
\]
\end{theorem}

\begin{proof}
The second formula is the root filter expression for $\rho^{-j}\Efun_{N,j}(\rho z)$.  The first is the case $j=0$.
\end{proof}

\begin{definition}\label{def:reduced_denominator}
Define the reduced denominator
\[
\Delta_N(w)
=
\sum_{n\geq0}(-1)^n\frac{w^n}{(Nn)!}.
\]
Then
\[
\Dfun_N(z)=\Delta_N(z^N).
\]
\end{definition}

\begin{corollary}[Symmetries]\label{cor:denominator_symmetries}
The denominator satisfies
\[
\Dfun_N(\omega z)=\Dfun_N(z)
\]
and
\[
\Dfun_N(\overline z)=\overline{\Dfun_N(z)}.
\]
Its nonzero zeros occur in orbits under multiplication by the $N$th roots of unity and under complex conjugation.
\end{corollary}

\begin{proof}
The first identity follows from the powers $z^{Nn}$.  The second follows from the real coefficients.  Both symmetries preserve the zero set and multiplicity.
\end{proof}

\begin{example}\label{ex:small_denominators}
The first denominators are
\[
\Dfun_1(z)=e^{-z},
\qquad
\Dfun_2(z)=\cos z.
\]
For $N=4$,
\[
\Dfun_4(z)
=
\cos\left(\frac z{\sqrt2}\right)
\cosh\left(\frac z{\sqrt2}\right).
\]
\end{example}

\section{Reciprocal powers and parameter polynomials}\label{sec:reciprocal_powers}

Since $\Dfun_N(0)=1$, its formal logarithm is defined at the origin.  On a simply connected zero free region containing the origin, the same logarithm is holomorphic.

\begin{definition}\label{def:reciprocal_family}
For $r\in\CC$, define
\[
\Sfun_N^{(r)}(z)
=
\Dfun_N(z)^{-r}
=
\exp\bigl(-r\log\Dfun_N(z)\bigr).
\]
Write
\[
\Sfun_N^{(r)}(z)
=
\sum_{n\geq0}P_{N,n}(r)\frac{z^{Nn}}{(Nn)!}.
\]
Define the logarithmic coefficients $\kappa_{N,m}$ by
\[
-\log\Delta_N(w)
=
\sum_{m\geq1}\kappa_{N,m}\frac{w^m}{m}.
\]
\end{definition}

The polynomial dependence on $r$ can be read directly from the formal logarithm.  We use the complete exponential Bell polynomials $B_n$ defined by
\[
\exp\left(\sum_{m\geq1}x_m\frac{t^m}{m!}\right)
=
\sum_{n\geq0}B_n(x_1,\ldots,x_n)\frac{t^n}{n!}.
\]

\begin{theorem}[Bell polynomial formula]\label{thm:reciprocal_bell}
For every $n\geq0$,
\[
P_{N,n}(r)
=
\frac{(Nn)!}{n!}
B_n\bigl(
 r\kappa_{N,1},
 r1!\kappa_{N,2},
 \ldots,
 r(n-1)!\kappa_{N,n}
\bigr).
\]
For $n\geq1$, this is a polynomial in $r$ of exact degree $n$.  Its leading coefficient is
\[
\frac{(Nn)!}{n!(N!)^n}.
\]
\end{theorem}

\begin{proof}
The identity
\[
\Delta_N(w)^{-r}
=
\exp\left(
 r\sum_{m\geq1}\kappa_{N,m}\frac{w^m}{m}
\right)
\]
fits the defining generating function for $B_n$ with
\[
x_m=r(m-1)!\kappa_{N,m}.
\]
The coefficient of $w^n$ is therefore the displayed Bell polynomial divided by $n!$.  Multiplication by $(Nn)!$ gives $P_{N,n}(r)$.  Since
\[
\Delta_N(w)=1-\frac{w}{N!}+O(w^2),
\]
we have $\kappa_{N,1}=1/N!$.  The term $x_1^n$ gives the stated leading coefficient, which is nonzero.
\end{proof}

The coefficients also satisfy a triangular recurrence that does not require the logarithmic coefficients to be known first.

\begin{theorem}[Parameter recurrence]\label{thm:parameter_recurrence}
We have $P_{N,0}(r)=1$.  For $n\geq1$,
\[
\begin{aligned}
P_{N,n}(r)
={}&-
\sum_{k=1}^{n-1}
(-1)^k
\binom{Nn-1}{Nk}
P_{N,n-k}(r)
\\
&-r
\sum_{k=1}^{n}
(-1)^k
\binom{Nn-1}{Nk-1}
P_{N,n-k}(r).
\end{aligned}
\]
\end{theorem}

\begin{proof}
Set $S=\Dfun_N^{-r}$.  Logarithmic differentiation gives
\[
\Dfun_NS'+r\Dfun_N'S=0.
\]
Take the coefficient of $z^{Nn-1}/(Nn-1)!$.  The term with $k=0$ in the first product is $P_{N,n}(r)$.  Every other term has a smaller block index.  Rearrangement gives the formula.
\end{proof}

\begin{theorem}[Reciprocal semigroup]\label{thm:reciprocal_semigroup}
For all $r,s\in\CC$,
\[
\Sfun_N^{(r)}(z)\Sfun_N^{(s)}(z)
=
\Sfun_N^{(r+s)}(z)
\]
as a formal identity and on every common zero free region with a fixed logarithm.  Consequently,
\[
P_{N,n}(r+s)
=
\sum_{q=0}^{n}
\binom{Nn}{Nq}
P_{N,q}(r)P_{N,n-q}(s).
\]
\end{theorem}

\begin{proof}
Both statements follow by adding the exponents of the same logarithm and then multiplying the block exponential generating series.
\end{proof}

\begin{corollary}[A sequence of binomial type]\label{cor:binomial_type}
Define
\[
A_{N,n}(r)
=
\frac{n!}{(Nn)!}P_{N,n}(r).
\]
Then
\[
A_{N,n}(r+s)
=
\sum_{q=0}^{n}
\binom nq
A_{N,q}(r)A_{N,n-q}(s).
\]
Thus $(A_{N,n})_{n\geq0}$ is a polynomial sequence of binomial type.
\end{corollary}

\begin{proof}
Multiply the coefficient identity in \cref{thm:reciprocal_semigroup} by $n!/(Nn)!$ and simplify the factorials.
\end{proof}

\begin{corollary}\label{cor:factor_r}
For $n\geq1$, the polynomial $P_{N,n}(r)$ is divisible by $r$.
\end{corollary}

\begin{proof}
At $r=0$, the reciprocal family is the constant function one.  Every positive block coefficient therefore vanishes.
\end{proof}

\begin{example}\label{ex:reciprocal_polynomials}
For $N=1$, one has $\Dfun_1(z)^{-r}=e^{rz}$ and $P_{1,n}(r)=r^n$.  For $N=2$,
\[
P_{2,1}(r)=r,
\]
\[
P_{2,2}(r)=r(3r+2),
\]
and
\[
P_{2,3}(r)=r(15r^2+30r+16).
\]
At $r=1$, these are the positive secant coefficients.
\end{example}

\section{Logarithmic derivatives and Riccati equations}\label{sec:tangent}

\begin{definition}\label{def:transfer_tangent}
Define the transfer tangent
\[
\Tfun_N(z)
=-\frac{\Dfun_N'(z)}{\Dfun_N(z)}.
\]
Its block expansion has the form
\[
\Tfun_N(z)
=
\sum_{n\geq0}T_N(n)
\frac{z^{Nn+N-1}}{(Nn+N-1)!}.
\]
\end{definition}

\begin{proposition}[Logarithmic coefficient formula]\label{prop:tangent_kappa}
The transfer tangent satisfies
\[
\Tfun_N(z)
=
N\sum_{m\geq1}\kappa_{N,m}z^{Nm-1}.
\]
Equivalently,
\[
T_N(n)
=
N(Nn+N-1)!\kappa_{N,n+1}.
\]
\end{proposition}

\begin{proof}
Since $\Dfun_N(z)=\Delta_N(z^N)$,
\[
-\frac{\Dfun_N'(z)}{\Dfun_N(z)}
=
-Nz^{N-1}\frac{\Delta_N'(z^N)}{\Delta_N(z^N)}.
\]
Differentiate the defining series for $-\log\Delta_N$ and compare coefficients.
\end{proof}

\begin{theorem}[Tangent recurrence]\label{thm:tangent_recurrence}
The coefficients satisfy $T_N(0)=1$ and
\[
T_N(n)
=
(-1)^n
-
\sum_{k=1}^{n}
(-1)^k
\binom{Nn+N-1}{Nk}
T_N(n-k).
\]
\end{theorem}

\begin{proof}
The defining equation is
\[
\Dfun_N\Tfun_N=-\Dfun_N'.
\]
Extract the coefficient of $z^{Nn+N-1}/(Nn+N-1)!$.
\end{proof}

\begin{corollary}[Derivative of the reciprocal family]\label{cor:reciprocal_derivative}
For every $r\in\CC$,
\[
\frac{d}{dz}\Sfun_N^{(r)}(z)
=
r\Tfun_N(z)\Sfun_N^{(r)}(z).
\]
\end{corollary}

\begin{proof}
Differentiate $\Dfun_N^{-r}$ on a zero free region.  The formal identity follows from the same logarithmic calculation.
\end{proof}

For a function $g$, define the complete differential Bell polynomial $Y_N$ by
\[
\frac{d^N}{dz^N}e^{g(z)}
=
e^{g(z)}Y_N\bigl(g',g'',\ldots,g^{(N)}\bigr).
\]

\begin{theorem}[Higher Riccati equation]\label{thm:higher_riccati}
On every zero free region of $\Dfun_N$,
\[
Y_N\bigl(
-\Tfun_N,
-\Tfun_N',
\ldots,
-\Tfun_N^{(N-1)}
\bigr)
=-1.
\]
\end{theorem}

\begin{proof}
Write $\Dfun_N=e^g$ locally.  Then $g'=-\Tfun_N$.  Divide the equation $\Dfun_N^{(N)}=-\Dfun_N$ by $\Dfun_N$ and use the defining identity for $Y_N$.
\end{proof}

\begin{example}\label{ex:ordinary_tangent}
For $N=1$, $\Tfun_1=1$.  For $N=2$, $\Dfun_2=\cos z$ and $\Tfun_2=\tan z$.  The Riccati equation becomes
\[
\Tfun_2'=1+\Tfun_2^2.
\]
\end{example}

\section{Congruential reciprocal functions and Appell systems}\label{sec:congruential}

The section $\Efun_{N,j}$ has a forced zero of order $j$ at the origin.  Removing it leaves a unit series in $z^N$.

\begin{definition}\label{def:normalized_section}
For $0\leq j<N$, define
\[
\Phi_{N,j}(w)
=
j!\sum_{n\geq0}\frac{w^n}{(Nn+j)!}.
\]
Then
\[
\Efun_{N,j}(z)
=
\frac{z^j}{j!}\Phi_{N,j}(z^N).
\]
Define the congruential reciprocal function
\[
\Cfun_{N,j}(z)
=
\Phi_{N,j}(z^N)^{-1}
=
\frac{z^j/j!}{\Efun_{N,j}(z)}.
\]
Write
\[
\Cfun_{N,j}(z)
=
\sum_{n\geq0}C_{N,j}(n)\frac{z^{Nn}}{(Nn)!}.
\]
\end{definition}

\begin{theorem}[Triangular inverse system]\label{thm:triangular_inverse}
The coefficients are determined by $C_{N,j}(0)=1$ and
\[
\sum_{q=0}^{n}
\binom{Nn+j}{Nq}
C_{N,j}(q)
=
\delta_{n0}.
\]
For $n\geq1$,
\[
C_{N,j}(n)
=
-
\frac{
\displaystyle\sum_{q=0}^{n-1}
\binom{Nn+j}{Nq}C_{N,j}(q)
}{
\displaystyle\binom{Nn+j}{Nn}
}.
\]
\end{theorem}

\begin{proof}
Multiply the defining series for $\Cfun_{N,j}$ and $\Efun_{N,j}$.  Their product is $z^j/j!$.  At degree $Nn+j$, multiplication by $(Nn+j)!$ gives the stated binomial coefficient.  The diagonal coefficient is nonzero, so the system is triangular.
\end{proof}

\begin{definition}\label{def:congruential_power}
For $r\in\CC$, define formal coefficients $C_{N,j,n}(r)$ by
\[
\Cfun_{N,j}(z)^r
=
\sum_{n\geq0}C_{N,j,n}(r)\frac{z^{Nn}}{(Nn)!}.
\]
On a zero free region, the same definition uses a fixed logarithm of $\Cfun_{N,j}$.
\end{definition}

\begin{theorem}[Congruential reciprocal semigroup]\label{thm:congruential_semigroup}
For each $n\geq1$, the coefficient $C_{N,j,n}(r)$ is a polynomial in $r$ of exact degree $n$.  Its leading coefficient is
\[
\frac{(Nn)!}{n!}
\left(-\frac{j!}{(N+j)!}\right)^n.
\]
It satisfies
\[
C_{N,j,n}(r+s)
=
\sum_{q=0}^{n}
\binom{Nn}{Nq}
C_{N,j,q}(r)C_{N,j,n-q}(s).
\]
In particular,
\[
C_{N,j,n}(1)=C_{N,j}(n).
\]
\end{theorem}

\begin{proof}
The unit series $\Phi_{N,j}(w)^{-1}$ has a formal logarithm with zero constant term.  Its linear term is
\[
-\frac{j!}{(N+j)!}w.
\]
Exponentiation after multiplication by $r$ shows that the coefficient of $w^n$ is a polynomial of degree at most $n$.  The $n$th power of the linear logarithmic term gives the displayed nonzero leading coefficient, so the degree is exactly $n$.  Multiplication of the powers with parameters $r$ and $s$ gives the coefficient identity.
\end{proof}

\begin{definition}\label{def:block_appell}
Define polynomials $E_m^{(N,j,r)}(x)$ by
\[
\Cfun_{N,j}(z)^re^{xz}
=
\sum_{m\geq0}E_m^{(N,j,r)}(x)\frac{z^m}{m!}.
\]
\end{definition}

\begin{theorem}[Appell laws]\label{thm:appell_laws}
The polynomials satisfy
\[
\frac{d}{dx}E_m^{(N,j,r)}(x)
=
mE_{m-1}^{(N,j,r)}(x)
\]
and
\[
E_m^{(N,j,r)}(x+y)
=
\sum_{q=0}^{m}
\binom mq
E_q^{(N,j,r)}(x)y^{m-q}.
\]
\end{theorem}

\begin{proof}
Differentiate the generating function with respect to $x$ for the first identity.  Multiply it by $e^{yz}$ and compare coefficients for the second.
\end{proof}

\begin{proposition}[Explicit Appell polynomial]\label{prop:explicit_appell}
For every $m\geq0$,
\[
E_m^{(N,j,r)}(x)
=
\sum_{0\leq n\leq m/N}
\binom{m}{Nn}
C_{N,j,n}(r)x^{m-Nn}.
\]
\end{proposition}

\begin{proof}
Multiply the block series for $\Cfun_{N,j}^r$ by the exponential series and collect the coefficient of total degree $m$.
\end{proof}


\begin{corollary}[Appell symmetry]\label{cor:appell_symmetry}
Each polynomial $E_m^{(N,j,r)}$ is monic of degree $m$, and
\[
E_m^{(N,j,r)}(\omega x)
=
\omega^mE_m^{(N,j,r)}(x).
\]
\end{corollary}

\begin{proof}
The term $n=0$ in \cref{prop:explicit_appell} is $x^m$.  Every other exponent is $m-Nn$, so replacement of $x$ by $\omega x$ multiplies every term by $\omega^m$.
\end{proof}

\begin{example}\label{ex:euler_boundaries}
For $(N,j)=(2,0)$,
\[
\Cfun_{2,0}(z)=\sech z.
\]
For $(N,j)=(2,1)$,
\[
\Cfun_{2,1}(z)=\frac z{\sinh z}.
\]
The first case gives Euler polynomials after multiplication by $e^{xz}$.  The second belongs to the familiar Bernoulli type family.
\end{example}

\section{Scaled reciprocal coupling and theta polynomials}\label{sec:theta}

\begin{definition}\label{def:unsigned_denominator}
Set
\[
H_N(z)=\Efun_{N,0}(z)
=
\sum_{n\geq0}\frac{z^{Nn}}{(Nn)!}.
\]
For $u\in\CC$, define $\Thetafun_{N,n}(u)$ by
\[
H_N(uz)^{-1}H_N(z)^{-u}
=
\sum_{n\geq0}\Thetafun_{N,n}(u)\frac{z^{Nn}}{(Nn)!}.
\]
\end{definition}

\begin{theorem}[Theta reconstruction]\label{thm:theta_reconstruction}
For every $n\geq0$,
\[
\Thetafun_{N,n}(u)
=
(-1)^n
\sum_{q=0}^{n}
\binom{Nn}{Nq}
P_{N,q}(1)u^{Nq}P_{N,n-q}(u).
\]
Thus $\Thetafun_{N,n}$ is a polynomial of degree at most $Nn$.
\end{theorem}

\begin{proof}
Choose $\rho$ with $\rho^N=-1$.  Then
\[
H_N(z)=\Dfun_N(\rho z).
\]
The coefficient of $z^{Nn}$ in $H_N(z)^{-r}$ is therefore $(-1)^nP_{N,n}(r)/(Nn)!$.  Apply this to $H_N(uz)^{-1}$ and $H_N(z)^{-u}$, then multiply the two block series.
\end{proof}

\begin{proposition}[Trivial parameter factors]\label{prop:theta_factors}
For $n\geq1$, the polynomial $\Thetafun_{N,n}(u)$ is divisible by $u$.  If $N$ is even, it is also divisible by $u+1$.
\end{proposition}

\begin{proof}
At $u=0$, the defining product is one.  Hence every positive coefficient vanishes.  If $N$ is even, then $H_N(-z)=H_N(z)$.  At $u=-1$, the product is again one, so every positive coefficient vanishes there as well.
\end{proof}

\begin{example}\label{ex:theta_even}
For $N=2$,
\[
\Thetafun_{2,1}(u)=-u(u+1)
\]
and
\[
\Thetafun_{2,2}(u)
=
u(u+1)(5u^2+u+2).
\]
The two factors predicted by \cref{prop:theta_factors} are visible.
\end{example}

\section{Growth and the geometry of the large zeros}\label{sec:zero_geometry}

We now assume $N\geq2$.  Write
\[
\eta_k=e^{(2k+1)\pi i/N},
\qquad 0\leq k<N.
\]
These are the vertices of a regular polygon and
\[
\Dfun_N(z)=\frac1N\sum_{k=0}^{N-1}e^{\eta_kz}.
\]

\begin{standardfact}[Order from Taylor coefficients]\label{sf:coefficient_order}
If $f(w)=\sum_{n\geq0}a_nw^n$ is an entire function with infinitely many nonzero coefficients, then
\[
\ord f
=
\limsup_{n\to\infty}
\frac{n\log n}{\log(1/|a_n|)}.
\]
Canonical products for entire functions of order below one have genus zero.  We use these standard facts in the form given in \cite{Boas1954,Titchmarsh1939}.
\end{standardfact}

\begin{proposition}[Orders and exponential type]\label{prop:orders}
For $N\geq2$,
\[
\ord\Delta_N=\frac1N.
\]
The function $\Dfun_N$ has order one and exponential type one.
\end{proposition}

\begin{proof}
The coefficient of $w^n$ in $\Delta_N$ has modulus $1/(Nn)!$.  Stirling's formula and \cref{sf:coefficient_order} give order $1/N$.  Composition with $z^N$ gives order one for $\Dfun_N$.

The finite exponential representation gives
\[
|\Dfun_N(z)|\leq e^{|z|},
\]
so the type is at most one.  Choose a root $\eta_0$ and put $z=t\overline{\eta_0}$ with $t>0$.  The term $e^{\eta_0z}=e^t$ has strictly larger real exponent than every other term.  Hence
\[
\Dfun_N(t\overline{\eta_0})
=
\frac1Ne^t(1+o(1)).
\]
The type is therefore one.
\end{proof}

The zero strings are governed by two adjacent exponents.  Put
\[
a_N=\frac\pi N,
\qquad
s_N=\sin a_N,
\qquad
c_N=\cos a_N,
\]
and
\[
A_N=\frac\pi{s_N},
\qquad
 t_m=A_N\left(m+\frac12\right).
\]
For $N\geq3$, also put
\[
\gamma_N=c_N-\cos(3a_N)>0.
\]

\begin{lemma}[Two term dominance]\label{lem:two_term_dominance}
Assume $N\geq3$.  There are positive constants $Y$, $R$, and $C$ with the following properties.

If $z=x+iy$ lies in the sector $|\arg z|\leq a_N$, has $x\geq R$, and is a zero of $\Dfun_N$, then $|y|\leq Y$.

On the strip $x\geq R$ and $|y|\leq Y$, one has
\[
N e^{-e^{-ia_N}z}\Dfun_N(z)
=
1+e^{2is_Nz}+R_N(z),
\]
where
\[
|R_N(z)|+|R_N'(z)|
\leq
C e^{-\gamma_Nx}.
\]
\end{lemma}

\begin{proof}
The two exponents $e^{ia_N}$ and $e^{-ia_N}$ have the largest real parts on the positive ray.  Every other exponent has real part at most $\cos(3a_N)$ there.  If $|y|\leq Y$, division by $e^{e^{-ia_N}z}$ therefore gives the stated estimate, with the factor depending on $Y$ absorbed into $C$.

It remains to exclude zeros far from the central ray.  First remove a small angular neighborhood of the positive ray.  On the remaining compact set of directions in the sector, one vertex of the exponent polygon has a unique maximal real projection.  The corresponding exponential dominates the sum of the others when $|z|$ is large.

Inside the small angular neighborhood, suppose $y\geq Y$.  The exponential with exponent $e^{-ia_N}$ exceeds the one with exponent $e^{ia_N}$ by the factor $e^{2s_Ny}$.  Choose $Y$ so that this factor dominates the second term by a fixed margin.  The remaining exponentials are exponentially smaller in $x$.  Thus the largest term cannot be cancelled.  The case $y\leq-Y$ follows by conjugation.
\end{proof}

\begin{lemma}[Near cancellation of two exponentials]\label{lem:near_cancellation}
For every $M>0$, there are constants $\varepsilon_M>0$ and $C_M>0$ such that the following holds.  If
\[
|\Re w|\leq M
\]
and
\[
|1+e^w|\leq\varepsilon_M,
\]
then there is a unique integer $m$ with
\[
|w-(2m+1)\pi i|
\leq
C_M|1+e^w|.
\]
\end{lemma}

\begin{proof}
The function $1+e^w$ is $2\pi i$ periodic and has simple zeros at the odd multiples of $\pi i$.  Reduce the imaginary part to a compact interval of length $2\pi$.  Away from small disks about the zeros, the modulus has a positive minimum.  Inside each disk, the inverse function theorem gives the stated linear estimate.  Shrinking $\varepsilon_M$ makes the relevant zero unique.
\end{proof}

\begin{theorem}[Large zero strings]\label{thm:large_zero_strings}
Assume $N\geq3$.  There are constants $m_0$, $C$, and $\gamma>0$ and a sequence of positive real simple zeros $(\tau_m)_{m\geq m_0}$ of $\Dfun_N$ such that
\[
\tau_m
=
 t_m+O(e^{-\gamma t_m}).
\]
Every zero of sufficiently large modulus is of the form
\[
\omega^k\tau_m
\]
for a unique $0\leq k<N$ and a unique $m\geq m_0$.

The derivatives satisfy
\[
\Dfun_N'(\omega^k\tau_m)
=
\omega^{-k}\Dfun_N'(\tau_m)
\]
and
\[
|\Dfun_N'(\tau_m)|
=
\frac{2s_N}{N}e^{c_N\tau_m}
\left(1+O(e^{-\gamma\tau_m})\right).
\]
\end{theorem}

\begin{proof}
Work first in the sector $|\arg z|\leq a_N$.  By \cref{lem:two_term_dominance}, every sufficiently large zero lies in a fixed strip about the positive real axis.  Put
\[
G(z)=N e^{-e^{-ia_N}z}\Dfun_N(z).
\]
Then
\[
G(z)=1+e^{2is_Nz}+R_N(z)
\]
with an exponentially small remainder and derivative.

If $G(z)=0$, then
\[
|1+e^{2is_Nz}|
\leq
C e^{-\gamma_N\Re z}.
\]
The real part of $2is_Nz$ is bounded in the strip.  By \cref{lem:near_cancellation}, the point $z$ lies at distance $O(e^{-\gamma_N\Re z})$ from one of the real points
\[
 t_m=\frac{(2m+1)\pi}{2s_N}.
\]

Fix a large $m$.  On a circle
\[
|z-t_m|=K e^{-\gamma_Nt_m}
\]
with $K$ fixed and large, the function $1+e^{2is_Nz}$ has modulus at least a constant times $K e^{-\gamma_Nt_m}$.  The remainder is at most a smaller constant times $e^{-\gamma_Nt_m}$.  Rouch\'e's theorem gives exactly one zero of $G$ in the disk.  The count is with multiplicity, so this zero is simple.  The disk is invariant under complex conjugation and $\Dfun_N$ has real coefficients.  Uniqueness forces the zero to be real.  Denote it by $\tau_m$.

The localization argument shows that these disks contain every large zero in the fundamental sector.  Rotational invariance gives the zeros $\omega^k\tau_m$ and covers the plane.  The disks are disjoint for large $m$, which gives uniqueness of $k$ and $m$.

Differentiate the normalized expression.  At $z=\tau_m$,
\[
G'(\tau_m)
=
2is_Ne^{2is_N\tau_m}+R_N'(\tau_m)
=
-2is_N+O(e^{-\gamma\tau_m}).
\]
Since $G(\tau_m)=0$,
\[
\Dfun_N'(\tau_m)
=
\frac1N e^{e^{-ia_N}\tau_m}G'(\tau_m).
\]
Taking absolute values gives the derivative asymptotic.  Differentiating $\Dfun_N(\omega z)=\Dfun_N(z)$ gives the rotation formula.
\end{proof}

\begin{corollary}[The case $N=2$]\label{cor:case_two_zeros}
For $N=2$,
\[
\Dfun_2(z)=\cos z.
\]
Its zeros are
\[
\pm\left(m+\frac12\right)\pi,
\qquad m\geq0,
\]
and all are simple.  They agree with the formula $t_m=A_2(m+1/2)$.
\end{corollary}

\begin{proof}
This is the standard zero set of the cosine function.
\end{proof}

\begin{corollary}[Eventual simplicity and exact rays]\label{cor:eventual_simplicity}
For every $N\geq2$, all but finitely many zeros of $\Dfun_N$ are simple.  Every sufficiently large zero lies on one of the rays
\[
\omega^k(0,\infty),
\qquad 0\leq k<N.
\]
\end{corollary}

\begin{proof}
Use \cref{thm:large_zero_strings} for $N\geq3$ and \cref{cor:case_two_zeros} for $N=2$.
\end{proof}


\begin{example}[Exact fourth order strings]\label{ex:fourth_order_zeros}
For $N=4$,
\[
\Dfun_4(z)
=
\cos\left(\frac z{\sqrt2}\right)
\cosh\left(\frac z{\sqrt2}\right).
\]
The positive zeros are exactly
\[
\tau_m
=
\sqrt2\pi\left(m+\frac12\right)
=
A_4\left(m+\frac12\right).
\]
Rotation by the fourth roots of unity gives the real and imaginary zero strings.  In this case the error term in \cref{thm:large_zero_strings} is identically zero.
\end{example}

\begin{corollary}[Reduced zeros]\label{cor:reduced_zeros}
Let $\Lambda_N$ be the nonzero zero multiset of $\Delta_N$.  All but finitely many members of $\Lambda_N$ are positive and simple.  They can be indexed as
\[
\lambda_m=\tau_m^N
\]
for $m\geq m_0$, and
\[
\lambda_m
=
 t_m^N
+O\left(t_m^{N-1}e^{-\gamma t_m}\right).
\]
\end{corollary}

\begin{proof}
The identity $\Dfun_N(z)=\Delta_N(z^N)$ identifies each orbit of $N$ nonzero zeros of $\Dfun_N$ with one zero of $\Delta_N$.  The large orbit representatives are the positive zeros $\tau_m$.  Simplicity is preserved because
\[
\Dfun_N'(z)=Nz^{N-1}\Delta_N'(z^N)
\]
and $z\ne0$.  The asymptotic follows by expanding $\tau_m^N$ about $t_m$.
\end{proof}

\begin{theorem}[Zero counting asymptotics]\label{thm:zero_counting}
Let $n_{\Dfun_N}(R)$ count the zeros of $\Dfun_N$ in $|z|\leq R$, with multiplicity.  Let $n_{\Delta_N}(R)$ count the zeros of $\Delta_N$ in $|w|\leq R$.  Then
\[
n_{\Dfun_N}(R)
=
\frac{N\sin(\pi/N)}{\pi}R+O(1)
\]
and
\[
n_{\Delta_N}(R)
=
\frac{\sin(\pi/N)}{\pi}R^{1/N}+O(1).
\]
\end{theorem}

\begin{proof}
Each of the $N$ rays contains one large zero near
\[
 t_m=\frac\pi{\sin(\pi/N)}\left(m+\frac12\right).
\]
The number of such indices with $t_m\leq R$ is
\[
\frac{\sin(\pi/N)}\pi R+O(1).
\]
Multiplication by $N$ and addition of the finite exceptional set give the first formula.  The second follows from
\[
n_{\Dfun_N}(R^{1/N})=Nn_{\Delta_N}(R).
\]
\end{proof}

\section{Hadamard products and zero power sums}\label{sec:hadamard}

The reduction from $z$ to $w=z^N$ removes the rotational repetition of the zeros.

\begin{theorem}[Genus zero product]\label{thm:genus_zero_product}
For $N\geq2$, let $\Lambda_N$ be the nonzero zero multiset of $\Delta_N$.  Then
\[
\Delta_N(w)
=
\prod_{\lambda\in\Lambda_N}
\left(1-\frac w\lambda\right),
\]
where the product converges locally uniformly.  Consequently,
\[
\Dfun_N(z)
=
\prod_{\lambda\in\Lambda_N}
\left(1-\frac{z^N}{\lambda}\right).
\]
\end{theorem}

\begin{proof}
By \cref{prop:orders}, the order of $\Delta_N$ is $1/N<1$.  By \cref{cor:reduced_zeros}, it has infinitely many zeros.  Hadamard factorization therefore gives a canonical product of genus zero.  The only possible exponential factor is a nonzero constant because the order is below one.  Since $\Delta_N(0)=1$, that constant is one.  Substitution of $w=z^N$ gives the second formula.
\end{proof}


\begin{proposition}[Products for the normalized sections]\label{prop:normalized_section_product}
For every $0\leq j<N$, the entire function
\[
\Phi_{N,j}(w)
=
j!\sum_{n\geq0}\frac{w^n}{(Nn+j)!}
\]
has order $1/N$.  If $N\geq2$ and $\Lambda_{N,j}$ is its nonzero zero multiset, then
\[
\Phi_{N,j}(w)
=
\prod_{\mu\in\Lambda_{N,j}}
\left(1-\frac w\mu\right).
\]
For $m\geq1$, put
\[
\sigma_{N,j,m}
=
\sum_{\mu\in\Lambda_{N,j}}\mu^{-m}.
\]
Then
\[
C_{N,j,n}(r)
=
\frac{(Nn)!}{n!}
B_n\bigl(
 r\sigma_{N,j,1},
 r1!\sigma_{N,j,2},
 \ldots,
 r(n-1)!\sigma_{N,j,n}
\bigr).
\]
\end{proposition}

\begin{proof}
The coefficient formula for the order gives $1/N$.  The function is not a polynomial.  An entire function of order below one with only finitely many zeros would be a polynomial times a constant by Hadamard factorization, which is impossible here.  The genus zero product follows as in \cref{thm:genus_zero_product}.

Taking the logarithm gives
\[
-\log\Phi_{N,j}(w)
=
\sum_{m\geq1}\sigma_{N,j,m}\frac{w^m}{m}.
\]
Since $\Cfun_{N,j}(z)^r=\Phi_{N,j}(z^N)^{-r}$, the Bell polynomial calculation from \cref{thm:reciprocal_bell} gives the coefficient formula.
\end{proof}

\begin{definition}\label{def:zero_power_sums}
For $m\geq1$, define
\[
\sigma_{N,m}
=
\sum_{\lambda\in\Lambda_N}\lambda^{-m},
\]
with multiplicities.
\end{definition}

The series converges absolutely because the exponent of convergence of $\Lambda_N$ is $1/N<1$.

\begin{theorem}[Newton sums for the reduced zeros]\label{thm:newton_sums}
For $|w|$ below the least zero modulus,
\[
-\log\Delta_N(w)
=
\sum_{m\geq1}\sigma_{N,m}\frac{w^m}{m}.
\]
Thus
\[
\kappa_{N,m}=\sigma_{N,m}.
\]
If
\[
a_m=\frac{(-1)^m}{(Nm)!},
\qquad a_0=1,
\]
then
\[
\sigma_{N,m}
=
-ma_m-
\sum_{k=1}^{m-1}a_k\sigma_{N,m-k}.
\]
\end{theorem}

\begin{proof}
Take the logarithm of the product in \cref{thm:genus_zero_product}.  Absolute convergence in a disk below the first zero permits termwise expansion
\[
-\log\left(1-\frac w\lambda\right)
=
\sum_{m\geq1}\frac{w^m}{m\lambda^m}.
\]
This gives the logarithmic formula and identifies $\kappa_{N,m}$ with $\sigma_{N,m}$.

Differentiate the product identity in logarithmic form to obtain
\[
\Delta_N'(w)
=
-\Delta_N(w)
\sum_{m\geq1}\sigma_{N,m}w^{m-1}.
\]
Comparison of the coefficient of $w^{m-1}$ gives
\[
ma_m
=
-\sum_{k=0}^{m-1}a_k\sigma_{N,m-k}.
\]
Solving for $\sigma_{N,m}$ gives the recurrence.
\end{proof}



\begin{corollary}[Mittag Leffler form of the transfer tangent]\label{cor:mittag_leffler_tangent}
Away from the zeros of $\Dfun_N$,
\[
\Tfun_N(z)
=
Nz^{N-1}
\sum_{\lambda\in\Lambda_N}
\frac1{\lambda-z^N}.
\]
The series converges locally uniformly away from the poles.
\end{corollary}

\begin{proof}
Logarithmic differentiation of the genus zero product gives
\[
-\frac{\Delta_N'(w)}{\Delta_N(w)}
=
\sum_{\lambda\in\Lambda_N}\frac1{\lambda-w}.
\]
The series is locally uniform because $\sum_{\lambda}|\lambda|^{-1}<\infty$.  Substitute $w=z^N$ and use the chain rule.
\end{proof}

\begin{proposition}[Theta Bell polynomial formula]\label{prop:theta_bell}
Put
\[
\tau_{N,m}=(-1)^m\sigma_{N,m}.
\]
Then
\[
\Thetafun_{N,n}(u)
=
\frac{(Nn)!}{n!}
B_n\bigl(
 \tau_{N,1}(u^N+u),
 1!\tau_{N,2}(u^{2N}+u),
 \ldots,
 (n-1)!\tau_{N,n}(u^{nN}+u)
\bigr).
\]
\end{proposition}

\begin{proof}
The reduced unsigned denominator is
\[
\Psi_N(w)=\sum_{n\geq0}\frac{w^n}{(Nn)!}=\Delta_N(-w).
\]
By \cref{thm:newton_sums},
\[
-\log\Psi_N(w)
=
\sum_{m\geq1}\tau_{N,m}\frac{w^m}{m}.
\]
With $w=z^N$, the logarithm of the theta generating function is
\[
\sum_{m\geq1}
\tau_{N,m}(u^{Nm}+u)\frac{w^m}{m}.
\]
The complete Bell polynomial formula follows.
\end{proof}

\begin{corollary}[Tangent coefficients as zero sums]\label{cor:tangent_zero_sums}
For every $n\geq0$,
\[
T_N(n)
=
N(Nn+N-1)!\sigma_{N,n+1}.
\]
\end{corollary}

\begin{proof}
Combine \cref{prop:tangent_kappa,thm:newton_sums}.
\end{proof}

\begin{corollary}[Reciprocal coefficients from zero data]\label{cor:reciprocal_zero_data}
For every $n\geq0$,
\[
P_{N,n}(r)
=
\frac{(Nn)!}{n!}
B_n\bigl(
 r\sigma_{N,1},
 r1!\sigma_{N,2},
 \ldots,
 r(n-1)!\sigma_{N,n}
\bigr).
\]
\end{corollary}

\begin{proof}
Use \cref{thm:reciprocal_bell,thm:newton_sums}.
\end{proof}

\begin{example}[The cosine orbit sums]\label{ex:cosine_orbit_sums}
For $N=2$,
\[
\Delta_2(w)=\cos\sqrt w.
\]
Its zeros are
\[
\lambda_m=\pi^2\left(m+\frac12\right)^2,
\qquad m\geq0.
\]
Hence
\[
\sigma_{2,q}
=
\frac{2^{2q}-1}{\pi^{2q}}\zeta(2q).
\]
The identity in \cref{cor:tangent_zero_sums} recovers the classical tangent coefficients from these zero power sums.
\end{example}

\section{Coefficient asymptotics from the nearest zeros}\label{sec:coefficient_asymptotics}

For an entire function $f$, let $\Zcal(f)$ denote its multiset of nonzero zeros.  Let $r$ be a positive integer.  Then $\Dfun_N^{-r}$ is meromorphic and a simple zero of $\Dfun_N$ becomes a pole of order $r$.

\begin{theorem}[Full nearest circle expansion]\label{thm:nearest_circle_expansion}
Suppose the zeros of $\Dfun_N$ of least modulus $R$ are simple and form the finite set $\Zcal_R$.  Suppose the next zero modulus is $R_2>R$.  At each $\rho\in\Zcal_R$, write the principal part
\[
\Dfun_N(z)^{-r}
=
\sum_{\ell=1}^{r}
 a_{\rho,\ell}
\left(1-\frac z\rho\right)^{-\ell}
+
O(1).
\]
Then, for every $R_*$ with $R<R_*<R_2$ and $m=Nn$,
\[
P_{N,n}(r)
=
 m!
\sum_{\rho\in\Zcal_R}
\sum_{\ell=1}^{r}
 a_{\rho,\ell}
\binom{m+\ell-1}{\ell-1}
\rho^{-m}
+
O_{R_*}(m!R_*^{-m}).
\]
The leading Laurent coefficient is
\[
a_{\rho,r}
=
(-\rho\Dfun_N'(\rho))^{-r}.
\]
\end{theorem}

\begin{proof}
Subtract the complete principal parts at the zeros on $|z|=R$.  The remainder is holomorphic in every closed disk of radius $R_*<R_2$.  Cauchy's estimate gives an ordinary coefficient bound $O_{R_*}(R_*^{-m})$.  The coefficient of $z^m$ in
\[
\left(1-\frac z\rho\right)^{-\ell}
\]
is
\[
\binom{m+\ell-1}{\ell-1}\rho^{-m}.
\]
Multiplication by $m!$ converts the ordinary coefficient to $P_{N,n}(r)$.

Near a simple zero,
\[
\Dfun_N(z)
=
-\rho\Dfun_N'(\rho)
\left(1-\frac z\rho\right)
+O\left(\left(1-\frac z\rho\right)^2\right).
\]
Raising to the power $-r$ gives the formula for $a_{\rho,r}$.
\end{proof}

\begin{corollary}[Leading dominant zero term]\label{cor:leading_dominant_zero}
Under the hypotheses of \cref{thm:nearest_circle_expansion},
\[
\begin{aligned}
P_{N,n}(r)
={}&m!\binom{m+r-1}{r-1}
\sum_{\rho\in\Zcal_R}
\frac{\rho^{-m}}{(-\rho\Dfun_N'(\rho))^r}
\\
&+O\left(m!m^{r-2}R^{-m}\right)
+O_{R_*}\left(m!R_*^{-m}\right).
\end{aligned}
\]
For $r=1$, the first error term is absent.
\end{corollary}

\begin{proof}
The terms with $\ell<r$ in \cref{thm:nearest_circle_expansion} are bounded by a constant times $m^{r-2}R^{-m}$.  The term with $\ell=r$ gives the displayed main term.
\end{proof}

\begin{corollary}[Secant asymptotic]\label{cor:secant_asymptotic}
For $N=2$ and $r=1$,
\[
P_{2,n}(1)
\sim
2(2n)!\left(\frac2\pi\right)^{2n+1}.
\]
\end{corollary}

\begin{proof}
The nearest zeros of $\cos z$ are $\rho=\pm\pi/2$, both simple.  Their contributions in \cref{cor:leading_dominant_zero} are equal.
\end{proof}

\section{Zero zeta functions}\label{sec:zero_zeta}

The zero counting theorem gives the exact abscissae of absolute convergence.  The exponentially accurate zero locations give more.  They compare the zero zeta with a Hurwitz zeta function.

\begin{definition}\label{def:absolute_zero_zetas}
Define
\[
\mathfrak Z_N^{\mathrm{abs}}(s)
=
\sum_{\rho\in\Zcal(\Dfun_N)}|\rho|^{-s}
\]
and
\[
\mathfrak Z_N^{\mathrm{red}}(s)
=
\sum_{\lambda\in\Lambda_N}|\lambda|^{-s},
\]
with multiplicities.  Each term means $e^{-s\log|\rho|}$ or $e^{-s\log|\lambda|}$.
\end{definition}

\begin{theorem}[Sharp absolute convergence half planes]\label{thm:zeta_convergence}
The series $\mathfrak Z_N^{\mathrm{abs}}$ converges absolutely and locally uniformly when
\[
\Re s>1.
\]
It fails to converge absolutely when $\Re s\leq1$.  The series $\mathfrak Z_N^{\mathrm{red}}$ converges absolutely and locally uniformly when
\[
\Re s>\frac1N.
\]
It fails to converge absolutely when $\Re s\leq1/N$.  Both functions are holomorphic on their open half planes of absolute convergence.
\end{theorem}

\begin{proof}
By \cref{thm:zero_counting}, the zeros of $\Dfun_N$ have linear counting density, while the reduced zeros have counting order $R^{1/N}$.  Partial summation gives absolute convergence in the stated open half planes.  The positive terms obtained after taking absolute values diverge on the boundary and to its left.  Local uniform convergence follows by the Weierstrass test on compact subsets.
\end{proof}

\begin{theorem}[Meromorphic continuation of the symmetry grouped zeta]\label{thm:zeta_continuation}
Put
\[
A_N=\frac\pi{\sin(\pi/N)}.
\]
There is an entire function $H_N$ such that
\[
\mathfrak Z_N^{\mathrm{red}}(s)
=
A_N^{-Ns}\zeta\left(Ns,\frac12\right)+H_N(s)
\]
on $\Re s>1/N$.  The right hand side gives a meromorphic continuation to $\CC$ with one simple pole at $s=1/N$.  Its residue is
\[
\frac{\sin(\pi/N)}{N\pi}.
\]

There is an entire function $K_N$ such that
\[
\mathfrak Z_N^{\mathrm{abs}}(s)
=
N A_N^{-s}\zeta\left(s,\frac12\right)+K_N(s).
\]
Thus $\mathfrak Z_N^{\mathrm{abs}}$ continues meromorphically to $\CC$ with one simple pole at $s=1$, of residue
\[
\frac{N\sin(\pi/N)}\pi.
\]
\end{theorem}

\begin{proof}
By \cref{cor:reduced_zeros}, all but finitely many reduced zeros can be written
\[
\lambda_m
=
\left(A_N\left(m+\frac12\right)\right)^N(1+\varepsilon_m)
\]
with $\varepsilon_m$ exponentially small.  On every compact set of $s$ values,
\[
(1+\varepsilon_m)^{-s}-1
=
O(|\varepsilon_m|)
\]
uniformly for large $m$.  The difference
\[
\lambda_m^{-s}
-
A_N^{-Ns}\left(m+\frac12\right)^{-Ns}
\]
therefore has an exponentially convergent sum on every compact subset of $\CC$.  The finitely many exceptional zeros and the finitely many omitted model terms contribute entire functions.  This proves the first formula with an entire remainder.

The Hurwitz zeta function has a meromorphic continuation with one simple pole at its first argument equal to one and residue one.  Composition with $Ns$ gives the stated pole and residue.

Every nonzero zero $\lambda$ of $\Delta_N$ gives $N$ zeros of $\Dfun_N$, each of modulus $|\lambda|^{1/N}$.  Hence
\[
\mathfrak Z_N^{\mathrm{abs}}(s)
=
N\mathfrak Z_N^{\mathrm{red}}(s/N).
\]
Substitution gives the second formula.
\end{proof}

\begin{definition}\label{def:branch_orbit_zeta}
Choose a logarithm for each member of the finite exceptional part of $\Lambda_N$.  Use the real logarithm for the eventual positive sequence.  Define
\[
\mathfrak Z_N^{\mathrm{orb}}(s)
=
\sum_{\lambda\in\Lambda_N}e^{-s\Log\lambda}.
\]
\end{definition}

\begin{corollary}[Branch grouped continuation]\label{cor:branch_grouped_continuation}
The orbit zeta $\mathfrak Z_N^{\mathrm{orb}}$ continues meromorphically to $\CC$.  It has the same single pole and residue as $\mathfrak Z_N^{\mathrm{red}}$.  Changing the logarithms of the finite exceptional zeros changes the continuation by an entire function.
\end{corollary}

\begin{proof}
The tail consists of positive reduced zeros, so it has the same comparison with the Hurwitz zeta.  Only finitely many terms depend on the branch choices.
\end{proof}

\begin{definition}\label{def:derivative_weighted_zeta}
Let $q$ be a positive integer.  Choose logarithms of the simple zeros of $\Dfun_N$ with arguments in a fixed bounded interval.  Define
\[
\mathfrak Z_{N,q}^{\mathrm{der}}(s)
=
\sum_{\substack{\Dfun_N(\rho)=0\\ \Dfun_N'(\rho)\ne0}}
\frac{e^{-s\Log\rho}}{\Dfun_N'(\rho)^q}.
\]
\end{definition}

\begin{theorem}[Derivative weighted convergence]\label{thm:derivative_weighted_convergence}
If $N\geq3$ and $q\geq1$, the series $\mathfrak Z_{N,q}^{\mathrm{der}}(s)$ converges locally uniformly for every $s\in\CC$.  It therefore defines an entire function.

For $N=2$, the series converges absolutely on $\Re s>1$.
\end{theorem}

\begin{proof}
For $N\geq3$, \cref{thm:large_zero_strings} gives
\[
|\Dfun_N'(\rho)|^{-1}
\leq
C e^{-c_N|\rho|}
\]
for every sufficiently large simple zero.  On a compact set of $s$ values, the factor $e^{-s\Log\rho}$ grows at most polynomially in $|\rho|$.  Exponential decay therefore gives uniform convergence.  The finite exceptional set causes no difficulty.

For $N=2$, the derivative of the cosine has modulus one at every zero.  The assertion then follows from the linear zero spacing.
\end{proof}

\begin{remark}\label{rem:multiple_zeros}
The derivative weighted sum omits the finitely many zeros that might be multiple.  A treatment of those zeros requires the complete principal parts of $\Dfun_N^{-q}$ rather than the single derivative weight.
\end{remark}

\section{Holomorphic operator substitution}\label{sec:operator_substitution}

The scalar identities have a direct functional calculus form.  This section records only the basic consequence and does not develop an operator resolvent theory.

\begin{standardfact}[Holomorphic functional calculus]\label{sf:holomorphic_calculus}
Let $A$ be a bounded operator on a complex Banach space.  If $f$ is holomorphic on a neighborhood of $\Spec(A)$, then
\[
f(A)
=
\frac1{2\pi i}\int_\Gamma f(z)(zI-A)^{-1}\,dz
\]
for a contour $\Gamma$ surrounding the spectrum.  The calculus preserves sums, products, and composition.  We use the standard form from \cite{DunfordSchwartz1958}.
\end{standardfact}

\begin{proposition}[Operator reciprocal family]\label{prop:operator_reciprocal}
Let $A$ be bounded and suppose $\Spec(A)$ lies in a simply connected region on which $\Dfun_N$ has no zeros.  Fix a logarithm of $\Dfun_N$ there.  Then
\[
\Sfun_N^{(r)}(A)
=
\exp\bigl(-r\log\Dfun_N(A)\bigr)
\]
is well defined and
\[
\Sfun_N^{(r)}(A)\Sfun_N^{(s)}(A)
=
\Sfun_N^{(r+s)}(A).
\]
If $(P_i)$ is a finite complete family of bounded idempotents, then
\[
P_i f(A)P_j
=
\frac1{2\pi i}
\int_\Gamma
f(z)P_i(zI-A)^{-1}P_j\,dz
\]
for every holomorphic $f$ on a neighborhood of the spectrum.
\end{proposition}

\begin{proof}
The spectral mapping theorem makes $\Dfun_N(A)$ invertible.  The product identity follows from multiplicativity of the functional calculus and the fixed logarithm.  Applying $P_i$ and $P_j$ to the Cauchy integral formula gives the packet formula.
\end{proof}

\section{Further questions}\label{sec:further_questions}

The large zero theorem settles eventual simplicity, but it does not classify the finite exceptional set.  The strongest remaining zero question is whether every zero of $\Dfun_N$ is simple and whether every reduced zero of $\Delta_N$ is positive.  Both statements hold for $N=2$ and $N=4$.  The argument in this paper proves them outside a compact set for every $N$.

The approximation
\[
\tau_m
=
A_N\left(m+\frac12\right)
+O(e^{-\gamma m})
\]
can be expanded further by repeated substitution in the normalized two exponential equation.  A full expansion would refine the entire correction in the zero zeta formula and could lead to explicit special values.

The symmetry grouped zero zeta has a meromorphic continuation with one pole.  A functional equation is not obtained here.  The unsymmetrized branch zeta over all $N$ roots of each reduced zero also depends on how logarithms are assigned to the rays.  Any proposed functional equation must specify that choice.

The coefficient families in \cref{sec:reciprocal_powers,sec:congruential,sec:theta} have exact closure laws.  Their arithmetic behavior, including odd prime congruences and the separate dyadic weighting, belongs to a different problem.  It is not inferred from the complex analytic results proved here.

\section{Conclusion}\label{sec:conclusion}

Cyclic projection turns one exponential into a finite differential basis.  A signed cycle gives the denominator $\Dfun_N$.  Reciprocal powers of this denominator form a one parameter semigroup, and their normalized coefficients form a sequence of binomial type.  The logarithmic derivative satisfies a higher Riccati equation.  Normalized exponential sections produce congruential reciprocal functions and ordinary Appell systems.  A scaled reciprocal coupling produces the theta polynomials.

The change of variable $w=z^N$ removes the repeated rotational copies of the zeros.  The reduced denominator has order $1/N$ and a genus zero Hadamard product.  Its zero power sums are exactly the logarithmic coefficients that govern the transfer tangent and the reciprocal Bell polynomials.  Two adjacent exponentials determine every sufficiently large zero.  This gives eventual simplicity, the exact zero counting constant, derivative growth, and meromorphic continuation of the symmetry grouped zero zeta.

The common analytic chain is
\[
\begin{gathered}
\text{cyclic projection}
\longrightarrow
\text{finite differential basis}
\longrightarrow
\text{reciprocal powers}
\\
\longrightarrow
\text{logarithmic coefficients}
\longrightarrow
\text{zero power sums}
\longrightarrow
\text{zero zeta}.
\end{gathered}
\]
Each step is exact under the hypotheses stated in the corresponding theorem.

\section*{Use of language model assistance}
An OpenAI language model was used as an organizational and mathematical drafting aid.  Its contributions included suggestions for the reciprocal Bell polynomial formulation in \cref{sec:reciprocal_powers}, the reduced Hadamard product in \cref{sec:hadamard}, the dominant exponential proof in \cref{sec:zero_geometry}, and the Hurwitz zeta comparison in \cref{sec:zero_zeta}.  The author assumes full responsibility for checking every statement, proof, and citation in any submitted version.  The model is not an author.

\bibliographystyle{aomalpha}
\bibliography{PSC_Paper_IV_Annals_Human_Edition}

\end{document}
